\section{Introducción}

La digitalización rural implica la adopción de tecnologías digitales en contextos agrícolas remotos, transformando prácticas tradicionales en procesos eficientes y conectados. En Colombia, donde muchas áreas rurales carecen de infraestructura básica, la digitalización es un desafío pero también una oportunidad para el desarrollo.

El Portal Agro Comercial del Huila aborda esta necesidad mediante una plataforma diseñada para entornos rurales, con funcionalidades offline y soporte comunitario. Esto permite a los productores acceder a información, vender productos y gestionar fincas sin depender de conectividad constante.

\section{Marco teórico y trabajos relacionados}

Estudios sobre digitalización rural, como los de la CEPAL (2020), enfatizan la importancia de TIC para el desarrollo rural. Investigaciones en "Smart Villages" (2023) exploran cómo la digitalización transforma comunidades rurales en centros de innovación.

Trabajos relacionados incluyen análisis de brechas digitales en América Latina (IICA, 2025) y casos de éxito en digitalización agrícola (FAO, 2023). Estos subrayan la necesidad de enfoques inclusivos que consideren limitaciones de infraestructura.

\section{Metodología}

\subsection{Análisis de Necesidades Rurales}
Se realizó un mapeo de comunidades rurales en el Huila, identificando barreras de acceso a tecnología y oportunidades de digitalización.

\subsection{Desarrollo de Soluciones}
Se implementaron herramientas adaptadas, como aplicaciones móviles con modo offline y capacitaciones presenciales.

\subsection{Evaluación}
Pruebas en veredas remotas midieron la adopción y impactos en productividad.

\section{Implementación del Portal Agro Comercial del Huila}

El portal incluye módulos específicos para digitalización rural, como sincronización offline y acceso a datos climáticos.

\subsection{Conectividad y Acceso}
Alianzas para instalar Wi-Fi comunitario y distribución de dispositivos accesibles.

\subsection{Capacitación Digital}
Programas de formación adaptados a realidades rurales, con instructores locales.

\section{Resultados e Impactos}

La digitalización ha aumentado el acceso a mercados en un 40\% en comunidades participantes, con mejoras en gestión de fincas.

\subsection{Reducción de Brechas}
Más productores conectados, con disminución en aislamiento digital.

\subsection{Impactos Socioeconómicos}
Aumento en ingresos y fortalecimiento de comunidades.

\section{Discusión}

\subsection{Desafíos de Implementación}
Conectividad y resistencia cultural son obstáculos principales.

\subsection{Comparación con Literatura}
Resultados consistentes con estudios globales, destacando el potencial de plataformas integradas.

\section{Conclusiones y Visión Futura}

La digitalización rural es fundamental para el desarrollo agrícola colombiano. El portal acelera este proceso, creando comunidades más conectadas y prósperas.

\subsection{Trabajo Futuro}
Expandir infraestructura y integrar más tecnologías.

\section{Implementación en el proyecto Portal Agrocomercial del Huila}

El Portal Agrocomercial del Huila promueve la digitalización rural mediante funcionalidades offline, capacitaciones presenciales y alianzas para conectividad. Permite a productores en áreas remotas registrar productos, acceder a pronósticos climáticos y vender en línea, superando barreras de infraestructura. Además, el portal fomenta el uso de dispositivos móviles para gestión agrícola, contribuyendo a la inclusión digital y el desarrollo sostenible en el Huila.